\chapter{Spazio delle ipotesi e bias induttivo}
\label{cap:bias}

\begin{inbreve}
Inferire una funzione da un numero finito di esempi è un problema mal posto. Senza assunzioni sulla funzione
cercata, gli esempi non dicono nulla sui casi non visti. Il capitolo lo mostra in uno spazio di ipotesi
discreto, quello delle funzioni booleane. Dimostra poi che un learner senza bias non può generalizzare. Infine
distingue i due modi di introdurre un bias: vincolare il linguaggio delle ipotesi o vincolare la strategia di
ricerca.
\end{inbreve}

\section{Un problema mal posto}

Inferire funzioni generali da dati noti è un \termine{problema (inverso) mal posto}: la soluzione può violare
\textbf{esistenza}, \textbf{unicità} o \textbf{stabilità}. In linea di principio la soluzione non è unica, e con
dati finiti non ci si può aspettare di trovare quella esatta. Per questo si lavora con uno spazio delle ipotesi
\textbf{ristretto}, e lo studio di ogni modello ruota attorno a due domande:
\begin{itemize}
  \item \textbf{cosa si può rappresentare?} Una funzione non rappresentabile in $H$ non può nemmeno essere appresa;
  \item \textbf{cosa si può apprendere?} Domanda che viene dopo la prima.
\end{itemize}

Per costruire un modello e un algoritmo bisogna fare \textbf{assunzioni} sulla natura della funzione target.

\begin{definizione}[Bias induttivo]
Il \termine{bias induttivo} è l'insieme delle assunzioni fatte sulla funzione target per poter generalizzare.
Può riguardare
\begin{itemize}
  \item i \textbf{vincoli sul modello}, cioè su quali ipotesi $H$ può esprimere: \termine{language bias};
  \item i \textbf{vincoli o le preferenze dell'algoritmo}, cioè la strategia di ricerca in $H$: \termine{search
        bias};
  \item entrambi.
\end{itemize}
\end{definizione}

Il ruolo del bias si vede con chiarezza imparando un \textbf{concetto}, cioè una funzione booleana. Ad esempio
$h_{\mathit{gatto}}(\x)=1$ se l'animale descritto da $\x$ è un gatto, $0$ altrimenti.

\section{Imparare una funzione booleana}

Si vuole imparare una funzione booleana sconosciuta $d = f(x_1,x_2,x_3,x_4)$ dai 7 esempi della
Tabella~\ref{tab:booleana}.

\begin{table}[ht]
\centering\small
\caption{Sette esempi di una funzione booleana sconosciuta di quattro variabili.}
\label{tab:booleana}
\begin{tabular}{ccccc|c}
\toprule
Esempio & $x_1$ & $x_2$ & $x_3$ & $x_4$ & $d$ \\
\midrule
1 & 0 & 0 & 1 & 0 & 0\\
2 & 0 & 1 & 0 & 0 & 0\\
3 & 0 & 0 & 1 & 1 & 1\\
4 & 1 & 0 & 0 & 1 & 1\\
5 & 0 & 1 & 1 & 0 & 0\\
6 & 1 & 1 & 0 & 0 & 0\\
7 & 0 & 1 & 0 & 1 & 0\\
\bottomrule
\end{tabular}
\end{table}

\subsubsection{Quante sono le ipotesi possibili}
Con 4 input binari esistono $2^4=16$ istanze di input distinte. Una funzione booleana assegna a ciascuna un output
in $\{0,1\}$, quindi le funzioni possibili sono $2^{16}=65\,536$. Sapere quale sia quella corretta richiede di
vedere \textbf{tutte} le coppie input-output. Dopo 7 esempi le 9 istanze non viste possono avere output
qualsiasi: restano $2^9 = 512$ funzioni compatibili con i dati. In generale, con input e output binari e
dimensione dell'input $n$:
\begin{equation}\label{eq:numero-booleane}
\abs{H} = 2^{\#\text{istanze di input}} = 2^{2^n}.
\end{equation}

Un modello che usa tutto questo spazio è una \termine{lookup table}, un \ingl{rote learner}. Memorizza gli
esempi e classifica $\x$ solo se coincide con un esempio già visto; altrimenti non risponde.
\uline{Nessun bias induttivo, nessuna generalizzazione.}

\section{Restringere lo spazio: le regole congiuntive}

Si introduce un \textbf{language bias}: si considerano solo \termine{regole congiuntive}, cioè congiunzioni
(AND) di letterali, come $h_1 = x_2$, $h_2 = x_1\wedge x_2$, $h_3 = \mathit{true}$, $h_4 = \neg x_1\wedge x_2$.
Ad esempio $h_1$ significa ``se $x_2=1$ allora $h(\x)=1$, altrimenti $0$''.

\begin{proprieta}{Numero di regole congiuntive}
Con $n$ variabili binarie le ipotesi congiuntive semanticamente distinte sono $3^n + 1$.
\end{proprieta}
\begin{proof}
In una congiunzione, per ciascuna delle $n$ variabili ci sono tre possibilità: compare come $x_i$, compare come
$\neg x_i$, oppure non compare (``non importa''). Le scelte sono indipendenti, quindi si ottengono $3^n$
congiunzioni. Tutte le congiunzioni che contengono una coppia $x_i\wedge\neg x_i$ valgono sempre 0. Sono
equivalenti tra loro e alla funzione $\mathit{false}$, che va aggiunta come una sola ipotesi in più. Totale
$3^n+1$.
\end{proof}

Con $n=4$ si passa da $65\,536$ a $3^4+1 = 82$ ipotesi.

\begin{definizione}[Consistenza e version space]
Un'ipotesi $h$ è \termine{consistente} con il training set TR se $h(\x) = d(\x)$ per ogni esempio
$\langle\x, d(\x)\rangle\in TR$. Il \termine{version space} $VS_{H,TR}$ è il sottoinsieme delle ipotesi di $H$
consistenti con tutti gli esempi di training.
\end{definizione}

In questo spazio ridotto esistono algoritmi efficienti che trovano il version space senza enumerare tutte le
ipotesi. Ma l'assunzione congiuntiva è \textbf{troppo restrittiva}: se il concetto target non è in $H$ non può
essere rappresentato. ``Se $x_1=1$ \emph{oppure} $x_2=1$ allora 1'' non è una congiunzione.

\section{L'unbiased learner}

L'idea opposta è scegliere un $H$ capace di esprimere \textbf{ogni concetto insegnabile}. È l'insieme di tutti i
sottoinsiemi dello spazio degli input $X$, cioè l'insieme delle parti $\mathcal{P}(X)$ (disgiunzioni,
congiunzioni, negazioni). Con $n=10$ input binari, $\abs{X} = 2^{10} = 1024$ e
$\abs{\mathcal{P}(X)} = 2^{1024}\approx 10^{308}$ concetti. $H$ contiene sicuramente il concetto target. Che
cosa succede alla generalizzazione?

\begin{teorema}[Un learner senza bias non può generalizzare]
Con un \ingl{unbiased learner} le uniche istanze classificate in modo non ambiguo dal version space sono gli
esempi di training stessi. Si ottiene di nuovo una lookup table.
\end{teorema}
\begin{proof}
Sia $x_i$ un'istanza non presente nel training set e sia $h\in VS$.
\begin{enumerate}
  \item Si consideri la funzione $h'$ identica a $h$ ovunque tranne che in $x_i$, dove $h'(x_i)\neq h(x_i)$.
  \item Poiché $H$ contiene \emph{tutte} le funzioni su $X$, anche $h'\in H$.
  \item $h$ e $h'$ coincidono su tutti gli esempi di training ($x_i$ non è tra questi), quindi anche $h'$ è
        consistente: $h'\in VS$.
  \item La corrispondenza $h\mapsto h'$ accoppia ogni ipotesi del VS che classifica $x_i$ come positivo con una che
        lo classifica negativo, e viceversa.
\end{enumerate}
Quindi ogni istanza non vista è classificata 1 da esattamente metà delle ipotesi del VS e 0 dall'altra metà. Il
version space non può dare alcuna risposta su $x_i$.
\end{proof}

\begin{esempio}[Un solo input]
Con un input binario $x$, le funzioni possibili sono $H=\{x,\ \neg x,\ 0,\ 1\}$. Il TR contiene solo $x=0\mapsto
d=0$. Il version space è $\{x, 0\}$: entrambe valgono 0 in $x=0$. Sull'istanza di test $x=1$, però, $x$ risponde
1 e $0$ risponde 0. Non c'è modo di decidere, a meno di usare tutto $X$ come training set.
\end{esempio}

\begin{idea}[Futilità dell'apprendimento senza bias]
Un learner che non fa \textbf{alcuna assunzione a priori} sull'identità del concetto target \textbf{non ha alcuna
base razionale} per classificare istanze non viste. Dovrebbe vedere come esempio ogni singola istanza di $X$. Il
bias, di restrizione o di preferenza, non serve solo all'efficienza: \uline{è necessario per la
generalizzazione}. Il bias però non dice ancora \emph{quale} soluzione generalizza meglio.
\end{idea}

\section{Sistemi induttivi e deduttivi equivalenti}

Un sistema induttivo è un algoritmo di apprendimento che usa lo spazio $H$. Riceve gli esempi e una nuova
istanza, e produce una classificazione o ``non so''. Un sistema induttivo è \textbf{equivalente a un sistema
deduttivo}, cioè a un dimostratore di teoremi. Il dimostratore riceve gli stessi esempi e la stessa istanza,
\textbf{più il bias induttivo come assioma aggiuntivo} (\figrif{fig:induttivo-deduttivo}). Il bias induttivo è
quindi esattamente l'insieme di assunzioni che rende \emph{logicamente deducibile} la classificazione prodotta dal
learner.

\begin{figure}[ht]
\centering
\begin{tikzpicture}[node distance=0.4cm]
\node[font=\small, align=right, anchor=east] (e1) at (0,0.3) {esempi di training};
\node[font=\small, align=right, anchor=east] (n1) at (0,-0.3) {nuova istanza};
\node[blocco, minimum width=4.2cm, minimum height=1.3cm] (s1) at (3.2,0) {Algoritmo di apprendimento\\che usa lo spazio $H$};
\node[font=\small, align=left, anchor=west] (o1) at (6.2,0) {classificazione\\o ``non so''};
\draw[->] (e1.east) -- (s1.west|-e1.east); \draw[->] (n1.east) -- (s1.west|-n1.east); \draw[->] (s1) -- (o1);
\node[font=\footnotesize, grigio, anchor=south west] at (s1.north west) {sistema induttivo};
\begin{scope}[yshift=-2.4cm]
\node[font=\small, align=right, anchor=east] (e2) at (0,0.45) {esempi di training};
\node[font=\small, align=right, anchor=east] (n2) at (0,0) {nuova istanza};
\node[font=\small, align=right, anchor=east] (b2) at (0,-0.45) {\textbf{bias induttivo}};
\node[blocco, minimum width=4.2cm, minimum height=1.3cm] (s2) at (3.2,0) {Dimostratore\\di teoremi};
\node[font=\small, align=left, anchor=west] (o2) at (6.2,0) {classificazione\\o ``non so''};
\draw[->] (e2.east) -- (s2.west|-e2.east); \draw[->] (n2.east) -- (s2.west|-n2.east); \draw[->] (b2.east) -- (s2.west|-b2.east); \draw[->] (s2) -- (o2);
\node[font=\footnotesize, grigio, anchor=south west] at (s2.north west) {sistema deduttivo equivalente};
\end{scope}
\end{tikzpicture}
\caption{Il bias induttivo rende un sistema induttivo equivalente a un sistema deduttivo.}
\label{fig:induttivo-deduttivo}
\end{figure}

\section{Language bias o search bias?}

Nel ML moderno si preferisce in genere il \textbf{search bias}. Si usano modelli flessibili, con spazi di
ipotesi molto espressivi e capacità di approssimazione universale (reti neurali, alberi di decisione). Così si
evita il language bias e \textbf{non si esclude a priori} la funzione target. Il bias induttivo resta, ma è
affidato alla \textbf{strategia di ricerca}: una ricerca \textbf{incompleta} che privilegia certe soluzioni
rispetto ad altre.

\begin{puntichiave}
\begin{itemize}
  \item Con $n$ input binari le funzioni booleane sono $2^{2^n}$; le congiunzioni di letterali $3^n+1$.
  \item Il version space contiene le ipotesi consistenti con tutti gli esempi.
  \item Senza bias induttivo il version space non classifica nessuna istanza nuova: lookup table.
  \item Language bias = vincoli su $H$; search bias = preferenze della ricerca. Il ML moderno usa $H$ molto
        espressivi e search bias. \textbf{Ogni approccio di ML ha un bias induttivo}; il problema è
        caratterizzarlo per ciascun modello.
\end{itemize}
\end{puntichiave}

\begin{domande}
\begin{enumerate}
  \item Che cos'è il bias induttivo? Differenza tra language bias e search bias.
  \item Dimostrare che un \ingl{unbiased learner} non può generalizzare.
  \item Quante sono le ipotesi con $n$ input binari senza vincoli? E con sole regole congiuntive?
  \item Che cos'è il version space?
\end{enumerate}
\end{domande}
